% Lösungen Einheit 3 – Grundwert berechnen
\einheitenkopf[][3 Grundwert]{Einheit 3 von 4 · Grundwert berechnen}

\erg{34}{a) Prozentwert \quad b) Prozentsatz \quad c) Grundwert \quad d) Prozentwert \quad e) Grundwert}
\erg{35}{a) $20\,€$ \quad b) $40\,€$ \quad c) $50\,€$ \quad d) $100\,€$ \quad e) Der Teil bleibt $10\,€$; je kleiner der Anteil in Prozent, desto größer das Ganze.}
\erg{36}{a) $1\,\% = 4\,$kg, $100\,\% = 400\,$kg \quad b) $1\,\% = 0{,}40\,€$, $100\,\% = 40\,€$ \quad c) $1\,\% = 2{,}20\,€$, $100\,\% = 220\,€$}
\erg{37}{a) $10\,\% = 3$ Kinder, $100\,\% = 30$ Kinder \quad b) frei sind $75\,\%$; $75\,\% = 45$ Plätze, $25\,\% = 15$, $100\,\% = 60$ Plätze}
\erg{38}{a) Teil gesucht: $15\,$kg \quad b) Ganzes gesucht: $60\,€$ \quad c) Prozentsatz gesucht: $20\,\%$ \quad d) Ganzes gesucht: $1\,\% = 0{,}16\,€$, $100\,\% = 16\,€$ (nicht $25\,\%$ von $4\,€ = 1\,€$)}
\erg{39}{a) Tim rechnet $25\,\%$ von $15\,€$, aber $15\,€$ sind schon $25\,\%$; richtig $15\,€ : 25 \cdot 100 = 60\,€$ \quad b) Paul rechnet $60\,\%$ von 36, aber 36 Kinder sind schon $60\,\%$; richtig $36 : 60 \cdot 100 = 60$ Kinder (weniger Befragte als Ja-Sager kann nicht sein)}
\erg{40}{a) $100\,\%$ sind 100-mal so viel wie $1\,\%$. \quad b) Bei $25\,\%$ ist das Ganze größer (viermal so groß wie der Teil); bei $125\,\%$ ist das Ganze kleiner als der Teil.}
\erg{41}{a) $1\,\% = 0{,}9$ Kinder, $100\,\% = 90$ Kinder \quad b) Ja, zwei Busse haben 100 Plätze, und $100 > 90$.}
